% Exported from Math++. Compile with pdflatex.
\documentclass[11pt,a4paper]{article}
\usepackage[T1]{fontenc}
\usepackage[utf8]{inputenc}
\usepackage{lmodern}
\usepackage{amsmath,amssymb,amsthm,mathtools}
\usepackage{graphicx,booktabs,xcolor}
\usepackage[hidelinks]{hyperref}
\usepackage[margin=25mm]{geometry}
\theoremstyle{plain}
\newtheorem{theorem}{Theorem}
\newtheorem{lemma}[theorem]{Lemma}
\newtheorem{corollary}[theorem]{Corollary}
\newtheorem{proposition}[theorem]{Proposition}
\theoremstyle{definition}
\newtheorem{definition}[theorem]{Definition}
\newtheorem{example}[theorem]{Example}
\theoremstyle{remark}
\newtheorem*{remark}{Remark}
\setlength{\parindent}{0pt}\setlength{\parskip}{0.6em}
\title{Graphs in two dimensions}
\date{}
\begin{document}
\maketitle
A \texttt{Figure:} line draws a graph. List what to draw, separated by commas, and give a range for
the variable with \texttt{x \textbackslash{}in [a, b]}. Text after the first \texttt{. } is the caption.

\section*{Curves}

\begin{figure}[htbp]
\centering
\includegraphics[width=0.8\linewidth]{graphs-2d-fig1.png}
\caption{Sine and cosine on two full turns, with a legend.}
\end{figure}
Your own functions and constants are used as they are defined above the figure.

\begin{equation*}
  g \equiv 9.81
\end{equation*}
\begin{equation*}
  h(t) \coloneqq  20 t - \frac{g}{2} t^2
\end{equation*}
\begin{figure}[htbp]
\centering
\includegraphics[width=0.8\linewidth]{graphs-2d-fig2.png}
\caption{Height of a ball thrown up at $20$ m/s.}
\end{figure}
\section*{Sliders}

A variable declared with a range gets a slider, and the figure redraws as you drag it.

\begin{equation*}
  a \coloneqq  1, \; a \in [0, 3]
\end{equation*}
\begin{equation*}
  b \coloneqq  0, \; b \in [-2, 2]
\end{equation*}
\begin{figure}[htbp]
\centering
\includegraphics[width=0.8\linewidth]{graphs-2d-fig3.png}
\caption{Amplitude $a$ and shift $b$.}
\end{figure}
\section*{Asymptotes and jumps}

The curve is broken at an asymptote instead of joined by a vertical line.

\begin{figure}[htbp]
\centering
\includegraphics[width=0.8\linewidth]{graphs-2d-fig4.png}
\caption{Tangent.}
\end{figure}
\begin{figure}[htbp]
\centering
\includegraphics[width=0.8\linewidth]{graphs-2d-fig5.png}
\caption{Two reciprocals.}
\end{figure}
At a jump, an open dot marks the end that is not included and a closed dot the end that is.

\begin{figure}[htbp]
\centering
\includegraphics[width=0.8\linewidth]{graphs-2d-fig6.png}
\caption{A step.}
\end{figure}
\section*{Equations}

\begin{figure}[htbp]
\centering
\includegraphics[width=0.8\linewidth]{graphs-2d-fig7.png}
\caption{A circle (equal scales are used when you give no range).}
\end{figure}
\begin{figure}[htbp]
\centering
\includegraphics[width=0.8\linewidth]{graphs-2d-fig8.png}
\caption{An ellipse and a hyperbola.}
\end{figure}
\begin{figure}[htbp]
\centering
\includegraphics[width=0.8\linewidth]{graphs-2d-fig9.png}
\caption{A parabola opening sideways.}
\end{figure}
\section*{Parametric and polar}

\begin{figure}[htbp]
\centering
\includegraphics[width=0.8\linewidth]{graphs-2d-fig10.png}
\caption{A Lissajous curve.}
\end{figure}
\begin{figure}[htbp]
\centering
\includegraphics[width=0.8\linewidth]{graphs-2d-fig11.png}
\caption{A cardioid.}
\end{figure}
\begin{figure}[htbp]
\centering
\includegraphics[width=0.8\linewidth]{graphs-2d-fig12.png}
\caption{A three-petal rose.}
\end{figure}
\section*{Inequalities}

An inequality shades the region where it holds. A dashed edge means the edge itself is not included.
Where regions overlap, the colours blend.

\begin{figure}[htbp]
\centering
\includegraphics[width=0.8\linewidth]{graphs-2d-fig13.png}
\caption{Two inequalities.}
\end{figure}
\begin{figure}[htbp]
\centering
\includegraphics[width=0.8\linewidth]{graphs-2d-fig14.png}
\caption{A ring.}
\end{figure}
\section*{Points and level curves}

\begin{figure}[htbp]
\centering
\includegraphics[width=0.8\linewidth]{graphs-2d-fig15.png}
\caption{Two marked points on a parabola.}
\end{figure}
\begin{figure}[htbp]
\centering
\includegraphics[width=0.8\linewidth]{graphs-2d-fig16.png}
\caption{Level curves.}
\end{figure}
\end{document}
